\problemstart{AP-02-001}{Map the Floating-Point Lattice}{Difficulty 2/5}{Chapter 2 \textbullet\ numpy}{Measure the next representable number as a function of magnitude instead of treating machine numbers as dense reals.}

        \subsection{Mathematical statement}


For a floating format $\tau$ and nonnegative representable $x$, define the local
forward spacing
\[
  \operatorname{ulp}_{\tau}^{+}(x)
  =\operatorname{nextafter}_{\tau}(x,+\infty)-x.
\]
Within a normal binade $[2^e,2^{e+1})$, the spacing is piecewise constant and
scales as $2^{e-p+1}$ for precision $p$.  At zero it equals the smallest positive
subnormal number.


        \subsection{Code problem}


Implement \code{local\_spacing} for finite nonnegative values and one requested
NumPy floating dtype, given as a \code{type} or an \code{np.dtype} instance.
Cast the values first, then compute the forward neighbor
in the same dtype.  Return spacings in that dtype and reject integer dtypes,
negative values, NaN, and infinity with \code{ValueError}.


        \begin{contractbox}
        \begin{Code}
        def local_spacing(values: np.ndarray, *, dtype: np.dtype | type) -> np.ndarray:
        \end{Code}
        \end{contractbox}

        \subsection{Theory--engineering gap inventory}

        \begin{gapbox}
        \textbf{Explicit gap.} Mathematical analysis assumes arbitrarily close reals; a floating format is a magnitude-dependent discrete lattice with subnormal and normal regimes.

        \textbf{Implicit gap exposed by the result.} The log--log staircase and the zero-spacing observation show why one universal absolute tolerance cannot describe all magnitudes or dtypes.
        \end{gapbox}

        \subsection{Acceptance evidence}

        \begin{evidencebox}
        \begin{itemize}
          \item Every result equals an independently computed same-dtype \code{nextafter} difference.
  \item Spacing is positive, zero uses the smallest subnormal, and powers of two expose jumps.
  \item The benchmark draws three dtype curves and prints spacing ratios; interpretation replaces a fabricated pass threshold. (Trust-based: the test suite does not machine-check benchmark output.)
        \end{itemize}
        \end{evidencebox}

        \subsection{Constraints}

        \begin{itemize}
          \item Use NumPy's \code{nextafter}; do not manipulate IEEE bit fields.
  \item Support float16, float32, and float64.
  \item This is a measurement problem, not an implementation of floating-point arithmetic.
        \end{itemize}

        \begin{sourcebox}
        This is an atomic adaptation, not copied solution code. Nearest primary sources:
        \begin{itemize}
          \item \href{https://numpy.org/doc/stable/reference/generated/numpy.finfo.html}{NumPy: numpy.finfo}
  \item \href{https://numpy.org/doc/stable/reference/generated/numpy.nextafter.html}{NumPy: numpy.nextafter}
        \end{itemize}
        \end{sourcebox}
